\chapter{Reading an Ultrasound}

\section*{Patient focus}
This chapter explains the examination or clinical problem in accessible language. Patients should use it to prepare questions and understand the purpose, experience, benefits, and limitations of imaging; local instructions and the official report take priority.


\label{ch:read_ultrasound}

Ultrasound is a highly operator-dependent, real-time imaging modality. Unlike CT and MRI, which produce static cross-sectional slices, ultrasound relies on the transmission of high-frequency sound waves and the interpretation of their returning echoes. Reading an ultrasound requires an understanding of acoustic terminology, diagnostic artifacts, and Doppler flow principles.

---

\section*{1. Acoustic Terminology: Echogenicity}
Tissues are described on ultrasound based on how strongly they reflect sound waves back to the transducer:

\begin{itemize}
    \item \textbf{Anechoic (Black)}: Complete transmission of sound waves with zero reflections. Characteristic of simple fluids (e.g., urine in the bladder, bile in the gallbladder, simple cysts, blood).
    \item \textbf{Hypoechoic (Dark Grey)}: Weak reflections. Characteristic of solid organs (e.g., liver parenchyma, renal cortex), muscle, and many cellular soft tissue tumors.
    \item \textbf{Isoechoic (Intermediate Grey)}: Echogenicity equal to the surrounding tissue.
    \item \textbf{Hyperechoic (Bright White)}: Strong reflections. Characteristic of dense structures that absorb or reflect sound waves highly (e.g., bone, calcifications, gallstones, renal sinus fat, fibrous tissue, gas/air).
\end{itemize}

---

\section*{2. Diagnostic Acoustic Artifacts}
Acoustic artifacts are not errors; they are physical phenomena essential for making a medical diagnosis:

\subsection*{Posterior Acoustic Shadowing}
When sound waves hit a highly attenuating structure (like a calcified gallstone or bone), almost all sound is reflected or absorbed. No sound waves travel deeper than the object, resulting in a \textbf{sharp black shadow} directly behind it. 
\textbullet \textbf{Clinical relevance}: Confirms the presence of a true calcification or stone (e.g., cholelithiasis, nephrolithiasis) rather than a soft tissue polyp or tumor.

\subsection*{Posterior Acoustic Enhancement}
When sound waves pass through a low-attenuating fluid-filled structure (like a simple cyst), very little sound is absorbed. The sound waves emerging from the fluid are much stronger than those passing through the adjacent solid tissue, resulting in a \textbf{bright white band} directly behind the structure.
\textbullet \textbf{Clinical relevance}: Confirms that a well-defined anechoic mass is a fluid-filled \textbf{simple cyst} rather than a solid hypoechoic tumor.

---

\section*{3. Doppler Ultrasound: Analyzing Flow}
Doppler imaging utilizes the shift in sound wave frequency caused by moving red blood cells to determine blood flow direction and velocity:

\subsection*{Color Doppler}
Color Doppler overlays flow information onto the grayscale image. Remember the mnemonic \textbf{BART}:
\textbullet \textbf{Blue Away / Red Toward}: Red indicates blood flowing \emph{toward} the transducer probe; blue indicates blood flowing \emph{away} from it. (Note: Red does \emph{not} inherently mean artery, and blue does \emph{not} inherently mean vein).
\textbullet \textbf{Clinical relevance}: Useful for checking vascular patency, identifying turbulent flow, or verifying vascularity within a solid mass.

\subsection*{Spectral Doppler}
Displays flow velocities over time as a waveform:
\textbullet \textbf{Arterial Flow}: Pulsatile waveforms. High-resistance beds (e.g., external carotid, extremity arteries) show sharp systolic peaks and low diastolic flow. Low-resistance beds (e.g., internal carotid, renal arteries) show continuous forward diastolic flow.
\textbullet \textbf{Venous Flow}: Phasic, non-pulsatile waveforms that vary gently with respiration.

---

\section*{4. Systematic Reading: Key Clinical Scans}

\subsection*{Abdominal Ultrasound (The Gallbladder)}
When evaluating the gallbladder for acute cholecystitis, review these four parameters:
1. \textbf{Lumen}: Look for hyperechoic stones inside the anechoic lumen. Verify that they demonstrate posterior acoustic shadowing and roll with gravity.
2. \textbf{Wall Thickness}: Normal is \textbf{< 3 mm}. Thickening > 3 mm suggests inflammation or edema.
3. \textbf{Pericholecystic Fluid}: Look for a dark, anechoic crescent of fluid surrounding the gallbladder wall.
4. \textbf{Sonographic Murphy's Sign}: Pain elicited when the transducer probe is pressed directly over the visualized gallbladder.

\subsection*{Pelvic and Obstetric Ultrasound}
\textbullet \textbf{Gynecological}: Identify the uterus and measure the endometrial stripe thickness. Locate the ovaries, noting the presence of normal anechoic follicles.
\textbullet \textbf{Obstetric (Early Pregnancy)}: Identify the gestational sac within the uterus. Look for the yolk sac and the fetal pole. Assess the \textbf{fetal heart rate (FHR)} using M-mode (motion mode) to avoid thermal heating of the fetus with spectral Doppler. Measure the \textbf{crown-rump length (CRL)} to date the pregnancy.

\subsection*{Point-of-Care Ultrasound (POCUS): The eFAST Exam}
In trauma, the Extended Focused Assessment with Sonography for Trauma (eFAST) is performed to detect free fluid (blood) or air:
1. \textbf{Right Upper Quadrant (RUQ)}: Check \textbf{Morison's pouch} (hepato-renal recess) for an anechoic strip of free fluid.
2. \textbf{Left Upper Quadrant (LUQ)}: Check the spleno-renal interface and subdiaphragmatic space.
3. \textbf{Pelvis (Suprapubic)}: Check posterior to the bladder (rectovesical pouch in males; pouch of Douglas in females).
4. \textbf{Subxiphoid Cardiac View}: Check the pericardial space for fluid (pericardial tamponade).
5. \textbf{Thoracic Views}: Check the pleural line for "lung sliding." An absence of lung sliding indicates a \textbf{pneumothorax}.

---

\section*{5. Structured Reporting Example: Acute Cholecystitis}

\begin{testbox}[Indications:]
45-year-old female with acute right upper quadrant pain, positive clinical Murphy's sign, and fever.
\end{testbox}
\begin{refbox}[Findings:]
\textbf{Liver}: Normal size and echogenicity. No focal hepatic mass or intrahepatic biliary ductal dilatation. Portal vein is patent with hepatopetal flow.\\
\textbf{Gallbladder}: Multiple mobile, hyperechoic focus calculations are noted within the gallbladder lumen, the largest measuring 1.8 cm. These foci demonstrate posterior acoustic shadowing. The gallbladder wall is diffusely thickened, measuring 5.2 mm. An anechoic stripe of pericholecystic fluid is present. The sonographic Murphy's sign is positive. The common bile duct is within normal limits, measuring 4.0 mm.\\
\textbf{Pancreas \& Kidneys}: Visually assessed portions of the pancreas are normal. The right kidney is normal; no hydronephrosis.
\end{refbox}
\begin{pearlbox}[Impression:]
Cholelithiasis and findings of acute calculous cholecystitis.
\end{pearlbox}

\clearpage
